\documentclass[10pt]{article}
\usepackage[margin=20mm]{geometry}
\usepackage{amsmath,amssymb,fontspec,unicode-math,xurl,hyperref}
\IfFontExistsTF{DejaVu Serif}{\setmainfont{DejaVu Serif}}{\setmainfont{Latin Modern Roman}}
\setmathfont{latinmodern-math.otf}
\setlength{\parindent}{0pt}\setlength{\parskip}{6pt}\sloppy
\title{Technical Verification of Noisy SIGNSGD Is More Differentially Private Than You (Might) Think}
\author{AI-assisted technical review}
\date{}
\begin{document}
\maketitle
Abstract: The complete published paper and embedded appendices were inspected. The scalar binary tradeoff calculation and the Gaussian limit for dense antipodal inputs withstand independent checks. Theorem 2, however, extends the latter calculation to arbitrary bounded vectors without the required dispersion condition. A fixed one-coordinate experiment embedded in every dimension disproves that extension and prevents the claimed uniform reduction of the Gaussian differential privacy parameter over the full input ball. The convergence statements also omit expectations that their proofs require; expected stationarity bounds can be recovered. Reported accuracies were not reproduced, and their privacy calibration requires a separate valid accountant. Scientific adjudication remains pending.\par
\section*{Identity and scope}
This report reviews Richeng Jin and Huaiyu Dai, ICML 2025, PMLR 267, I25-43769 (0400), OpenReview `thCqMz1ZXw`. All 22 physical pages were read, including references and Appendices A–C. Printed and physical page numbers agree. All six figures and all six tables were actually inspected, as were image renderings of the relevant theorem and proof pages. The source is the frozen published PDF; its hash and physical-page extraction records accompany this report.\par
No author repository was identified in the supplied bounded source search. Separate OpenReview attachments were unavailable in source preparation; this does not establish that no code or supplement exists. No author code, model or dataset was run or imported. Reviewer-authored finite hypothesis tests, scalar calculations and aggregation probabilities were computed using only the Python standard library; scripts, outputs and command receipts are preserved. Numerical grids support the analytic arguments below but are not substituted for their quantifiers.\par
\section*{Supported findings}
Problem 1 (The universal vector Gaussian limit and resulting uniform amplification fail).  Theorem 2, p. 6, quantifies over any vectors \(x,x'\) satisfying \(\|x\|_2,\|x'\|_2\le C\) and asserts convergence of their sign-mechanism tradeoff to \(G_\mu\) as \(d\to\infty\), with the bound in Equation (19). For fixed \(C,\sigma>0\), that bound gives \(\mu\le\mu_*=(2C/\sigma)\sqrt{2/\pi}\). Remark 2, p. 7, uses this reduction to claim \(\pi/2\) more training steps at the same overall privacy guarantee.\par
Equation reference: eq:binary. Set \(a=C/\sigma\), \(q=\Phi(-a)\) and, in every dimension, choose \[x_d=Ce_1,\qquad x'_d=-Ce_1.\] The first output bit has probabilities \(1-q\) and \(q\) of being positive, respectively. Every remaining output bit is an independent fair coin under both hypotheses. Those common ancillary bits contain no testing information. Thus the exact Neyman–Pearson tradeoff, for every \(d\), is \[
f_d(\alpha)=
\begin{cases}
1-\dfrac{1-q}{q}\alpha,&0\le\alpha\le q,\\[3pt]
\dfrac{q}{1-q}(1-\alpha),&q\le\alpha\le1.
\end{cases}\] In particular \(f_d(q)=q\). But the claimed limiting GDP curve satisfies \[G_{\mu_*}(q)=\Phi(a-\mu_*)>\Phi(-a)=q,\] since \(\mu_*<2a\). For \(C=\sigma=1\), the actual value is \(0.158655254\), whereas the proposed Gaussian bound requires at least \(0.275664737\). A still smaller admissible \(\mu\) only raises the required curve. The discrepancy persists as \(d\to\infty\); this is a counterexample to the theorem’s stated conclusion, not merely an absent proof detail. The piecewise linear curve also cannot become an exactly Gaussian tradeoff by appending uninformative coordinates.\par
More precisely, under the theorem’s comparison of all pairs in the full radius-\(C\) input ball, the smallest uniform GDP parameter remains exactly \(2C/\sigma\) for every \(d\). The raw Gaussian mechanism and post-processing give that upper bound. Conversely, [eq:binary] (\#eq:binary) at \(\alpha=q\) requires \(\mu\ge2C/\sigma\). Scalar sign compression can improve the tradeoff at other test levels without improving this particular uniform GDP parameter. This observation preserves the valid information-loss result in Section 4.1 and does not deny every possible form of privacy amplification.\par
The proof on pp. 6 and 13–15 replaces arbitrary inputs by the dense antipodes \(x'=-x=(C/\sqrt d)\mathbf1\), then applies a central limit argument. Maximizing an extremal likelihood ratio does not establish ordering of the entire testing tradeoff for all input pairs. The Lyapunov step also needs a condition excluding a nonvanishing single-coordinate contribution to the total privacy-loss variance. The sparse sequence above violates exactly that condition. Equality of the two privacy-loss variances is valid for the symmetric antipodal pair used in the calculation, not an automatic property of every pair merely because a worst-case tradeoff is symmetric.\par
For the restricted dense antipodal sequence, an independent binomial calculation approaches \(G_{2\sqrt{2/\pi}}\) when \(C=\sigma=1\): the largest discrepancy over 99 interior test levels is about \(0.02290\), \(0.001781\) and \(0.0002041\) at \(d=10,100,1000\), respectively. Thus a useful restricted asymptotic theorem is plausible and its calculation is not being rejected. It must state an appropriate triangular-array dispersion condition and quantify approximation error before being used as a finite-dimensional privacy guarantee. Letting \(\sigma\to\infty\) as well makes both privacy curves tend to perfect privacy; that unscaled convergence alone does not establish the asserted nontrivial multiplicative gain.\par
This Major finding affects the claimed general vector amplification, the \(\pi/2\) composition argument and the privacy interpretation of the 500-versus-320-round comparison in Appendix C.2.2 and Table 3, pp. 21–22. The arithmetic \(500(2/\pi)=318.31\) is correct; its privacy premise is not supplied by Theorem 2. Repair requires either a restricted theorem with finite approximation control or a different valid worst-case accountant, followed by recalibration of the comparisons. Specify the adjacency relation, per-example clipping, mini-batch averaging, sampling and composition. Record-level sensitivity can differ from the full-input-ball comparison above. Without the actual noise scales and accountant, this review does not claim that a particular trained release violated its advertised privacy budget, or that the displayed accuracies are numerically false.\par
Problem 2 (The convergence bounds require expectation and stationarity qualifications).  Theorems 3 and 4, Equations (33)–(34), pp. 7–8, bound empirical averages of gradient norms at random iterates without an expectation or a probability qualifier. Their proofs, especially Equations (72)–(74) on pp. 16–17 and (82)–(84) on p. 19, start with expected descent and drop the expectation around those gradient norms. Expected descent does not imply the displayed pathwise bound.\par
For example, let \(F(w)=\sqrt{1+w^2}\), \(F^*=1\), \(w_0=1\), \(L=c=d=M=1\), and use its exact gradient. Assumptions 1–4 hold, including zero stochastic-gradient variance. For any finite \(T\) and \(\sigma>0\), the event that every transmitted noisy sign is negative has strictly positive probability. On that event \(w_t=1+t\eta\), so every squared gradient is at least \(1/2\). With \(\eta=T^{-1/2}\) and \(\sigma=T^{1/4}\), the stated right-hand order tends to zero. It cannot be an unconditional pathwise inequality. For Scheme II, take \(M=T\) identical workers with independent noises; the all-negative-message event still has positive probability, while the stated right-hand order also tends to zero and the \(\ell_1\) gradient remains at least \(1/\sqrt2\). Neither example contradicts an expectation bound: the exceptional event becomes very unlikely.\par
Equation reference: eq:expectedI, eq:expectedII. A direct repair is available. Write \(k=\sqrt{2/\pi}/\sigma\) and interpret the unbiasedness and independence assumptions conditionally on the iterate. For \(|g_i|\le c\), \[\left|\mathbb E[\operatorname{sign}(g_i+N)]-kg_i\right|
\le \frac{k|g_i|^3}{6\sigma^2},\qquad N\sim N(0,\sigma^2).\] This follows by integrating \(1-e^{-u^2/2}\le u^2/2\), without an asymptotic Taylor remainder. If \(A_t\) is the average of the transmitted bits, then \(\mathbb E[A_t\mid w_t]=k\nabla F(w_t)+e_t\) with \(\|e_t\|_2^2\le dk^2c^6/(36\sigma^4)\) and \(\|A_t\|_2^2\le d\). Smoothness and Young’s inequality yield \[
\frac1T\sum_{t=0}^{T-1}\mathbb E\|\nabla F(w_t)\|_2^2
\le \frac{2(F(w_0)-F^*)/(\eta T)+Ld\eta}{k}
       +\frac{dc^6}{36\sigma^4}.\] At the stated learning rate this has the leading terms in Theorem 3, with the missing expectation restored and the dimension dependence made explicit. For independent workers, each coordinate of \(A_t\) has conditional variance at most \(1/M\). Bounding the sign-decision error by its mean absolute deviation similarly gives for Scheme II \[
\frac1T\sum_{t=0}^{T-1}\mathbb E\|\nabla F(w_t)\|_1
\le\frac{F(w_0)-F^*}{\eta T}+\frac{Ld\eta}{2}
       +\frac{\sigma d\sqrt{2\pi}}{\sqrt M}+\frac{dc^3}{3\sigma^2}.\] This argument avoids the variance-notation inconsistency between Assumption 3, which bounds variance by \(\sigma_{L,i}\), and Equation (80), which takes a square root of \(\sigma_{L,i}^2\). A proof following the latter route instead must distinguish variance from standard deviation. These repairs show why the omission is classified Minor rather than treated as a refutation of expected convergence.\par
The conclusions must concern expected stationarity, not proximity to a local optimum as stated in Remark 3: smooth nonconvex functions may have stationary saddles. Remark 4 also needs a lower growth bound on the number of workers, such as \(M=\Theta(T)\), to turn \(\sigma/\sqrt M\) with \(\sigma=O(\sqrt T/\mu_T)\) into \(O(1/\mu_T)\); the printed \(M=O(T)\) alone allows \(M=1\). A high-probability claim would require a separate concentration argument and confidence parameter. The paper explicitly excludes the effect of gradient clipping from this convergence analysis; the repaired bounds retain that restriction.\par
\section*{Other checks and limits of interpretation}
Theorem 1’s two-point likelihood test and the scalar Gaussian/logistic matching were checked directly. With \(|x|\le c\), the logistic scale \(s=c/\log(\Phi(c/\sigma)/\Phi(-c/\sigma))\) matches the extreme binary probabilities, hence the scalar tradeoff. The leading terms \(|x|^3/(6\sigma^2)\) and \(|x|^3/(12s^2)\) in Equations (14)–(17) describe bias of the scaled mean bit. They are not mean squared estimation errors. At \(c=\sigma=1\) and \(x=0\), both biases are zero, while the single-bit squared errors are \(\pi/2\simeq1.57080\) and \(4s^2\simeq1.43724\). The paper itself acknowledges the Gaussian variance disadvantage and that its aggregation method does not always outperform the logistic one, especially in Appendix C.2.1. Its narrower heterogeneous-data claim should retain those qualifications.\par
Independent convolution of the 100 Bernoulli messages in Figure 3 confirms a nonmonotone aggregation-error phenomenon under a specified tie rule. At \(\sigma=12\), the Gaussian mechanism has strictly negative-majority probability \(0.42577\) and tie probability \(0.07957\), giving \(0.46556\) with random tie-breaking. Counting a zero aggregate as wrong gives \(0.50535\) instead. The visible curve is compatible with the former value, but the displayed \(\operatorname{sign}\) formula does not specify tie handling. This needs explicit implementation documentation for even worker counts. It is not evidence that the plotted curve is fabricated. The independently computed curves were not substituted for the unavailable historical experiment.\par
The qualitative displays in Figures 4–6 and the six tables were checked for the stated trends. Similar test accuracies and some performance loss from majority vote are visible; those observations cannot establish equal privacy by themselves. The 32-to-1 payload calculation is correct for one binary gradient coordinate versus one 32-bit value, before framing, model broadcasts, aggregation representation and other communication costs. It should not be interpreted as a measured 32-fold end-to-end speedup. Server-side aggregation also cannot retroactively reduce the information already seen by a server in the stated local-privacy threat model.\par
Reported experiments specify datasets, worker counts, Dirichlet heterogeneity, clipping thresholds, a learning-rate grid, 500 rounds and ten repeats. Actual configurations, seeds, raw runs, checkpoint selection, the meaning of the reported dispersion, calibrated noise values and historical source were not available for reproduction. The Gaussian-input illustration in Appendix C is an aggregation example; its unbounded normal inputs do not themselves certify the bounded-input privacy hypothesis. No claim of complete empirical reproduction or benchmark equivalence is made.\par
\section*{Repair status}
One Major finding and one Minor finding are supported. The sparse counterexample is exact; finite enumeration checks it without a dimension-limit approximation. The restricted dense limit and corrected expected stationarity bounds were independently checked, so the report distinguishes surviving results from the failed generalization. Fresh independent acceptance rederived both the exact privacy counterexample and the surviving expected convergence bounds, including the Minor finding. Scientific adjudication remains pending; human reviews remain empty.\par
\end{document}
